% !TeX program = xelatex
\documentclass[11pt,a4paper]{article}
\usepackage[margin=23mm]{geometry}
\usepackage{fontspec}
\setmainfont{Latin Modern Roman}
\usepackage{amsmath,amssymb,mathtools}
\usepackage{xcolor,enumitem,booktabs,longtable,array,xurl,hyperref}
\definecolor{muted}{HTML}{53636C}
\hypersetup{colorlinks=true,urlcolor=blue,linkcolor=blue}
\setlength{\parindent}{0pt}
\setlength{\parskip}{5pt}
\newcommand{\R}{\mathbb R}
\newcommand{\E}{\mathbb E}
\newcommand{\N}{\mathbb N}
\newcommand{\ind}{\mathbf 1}
\newcommand{\Var}{\operatorname{Var}}
\newcommand{\Cov}{\operatorname{Cov}}
\newcommand{\supp}{\operatorname{supp}}
\DeclareMathOperator*{\argmin}{arg\,min}
\DeclareMathOperator*{\argmax}{arg\,max}
\newcommand{\entrymeta}[3]{{\sffamily\small\color{muted}\textbf{\texttt{#1}}\quad|\quad #2\par #3\par}}
\newcommand{\Problemref}[1]{\texttt{#1}}
\newcommand{\JELref}[2]{\texttt{#2} (JEL \texttt{#1})}
\begin{document}
\section*{Observational interference with observed and latent network confounding}
\textbf{Problem ID:} \texttt{C31-124}\quad \textbf{JEL:} \texttt{C31}\par
% BEGIN ECONOMICS STATEMENT
\label{problem:OP-0491}
\label{beyond:CI-12}
\entrymeta{CI-12}{Econometric theory and causal inference}{Research agenda formalization}
\subsection*{Definitions and assumptions}
Observe $n$ independent clusters of $N$ units, recording pretreatment network $G$, all observed covariates $X$, assignment vector $W\in\{0,1\}^N$, and outcomes $Y_i=Y_i(W)\in[0,1]$. Let $U$ contain the full latent potential-outcome table. In the single-cluster notation, $P$ denotes a complete-data law of $(G,X,U,W)$, with observed outcomes obtained by consistency. Define $q(w\mid G,X)=\Pr_P(W=w\mid G,X)$ and $q(w\mid G,X,U)=\Pr_P(W=w\mid G,X,U)$ for every $w\in\{0,1\}^N$; both are conditional probability mass functions under this same law. Require both $q(w\mid G,X)>0$ and $q(w\mid G,X,U)>0$ almost surely. Impose the sensitivity restriction, for every pair $w,w'$,
\[
\Gamma^{-1}\leq
\frac{q(w\mid G,X,U)/q(w'\mid G,X,U)}
     {q(w\mid G,X)/q(w'\mid G,X)}
\leq\Gamma
\]
almost surely, where $\Gamma\geq1$ is specified. Additional graph-local response and observed adjustment restrictions must be explicitly declared. A specified policy $h\in\{g,g'\}$ is a measurable probability kernel $h(w\mid G,X)$ on $\{0,1\}^N$, with $h(w\mid G,X)\geq0$ and $\sum_w h(w\mid G,X)=1$. Any restriction to a smaller pretreatment information set is part of the policy specification. Construct $W^h$ using an auxiliary random variable independent of $(G,X,U,W)$ so that, conditionally on $(G,X,U)$, its assignment probabilities are $h(w\mid G,X)$. Thus policy randomization is independent of the potential table conditionally on $(G,X)$; unconditional independence of $W^h$ and $U$ is not imposed. Bounded density ratios mean that, for the law and cluster size under consideration, there is a finite bound $K$ such that $h(w\mid G,X)/q(w\mid G,X)\leq K$ almost surely for every $w$ and both policies. Any uniform bound over laws or cluster sizes must be declared separately. Target
$\theta(P)=N^{-1}\sum_i\E_P[Y_i(W^g)-Y_i(W^{g'})]$.
For declared additional restrictions $\mathcal R$, let $\mathcal M_{\Gamma,\mathcal R}$ be the class of complete-data laws satisfying consistency, positivity, the displayed sensitivity restriction, the policy-ratio requirement, and $\mathcal R$. For an observed law $Q$ of $(G,X,W,Y)$, sharpness concerns exactly
\[
 \Theta(Q;\mathcal M_{\Gamma,\mathcal R})
 =\{\theta(P):P\in\mathcal M_{\Gamma,\mathcal R},\ P^{\rm obs}=Q\}.
\]
An empty set means that the observed law is incompatible with the declared class. This single-cluster notation does not impose identical distributions on the sampled clusters or specify how heterogeneous sampled clusters represent the target law.
\subsection*{Formal question}
Characterize when the observed law and the declared network restrictions identify $\theta$, and otherwise derive the sharp identified set over the sensitivity class. Determine which own-unit and neighbor covariates form sufficient adjustment sets when $\Gamma=1$, and which exposure reductions preserve the policy target. For bounded-degree response classes, obtain optimal estimation rates and honest confidence sets for the identified set, allowing $N$ to grow with $n$.
The network and causal restrictions are inputs to this family of characterization problems: each instance must declare the graph type, neighborhoods, permitted potential-response dependence, and any causal restrictions used for adjustment. A sufficient adjustment claim must exhibit a functional of the observed law using the retained covariates and show that it equals the policy contrast throughout the declared class; full-vector exchangeability given $(G,X)$ alone does not justify dropping covariates. An exposure reduction must specify its mappings and reduced observable information, and exhibit a functional of that information equal to the original contrast throughout the declared class for the specified policy pair. Preservation uniformly over a policy class is a stronger claim and must name that class. Bounded degree alone supplies neither a statistical regularity class nor independent replication within a cluster. Any rate or confidence-set claim must additionally declare the sampling experiment with $N=N_n$, its model classes and quantitative overlap bounds, nuisance information and dependence restrictions, the set-estimation loss, and the set-valued coverage and excess-size criteria; the scalar coverage convention below does not select those criteria.
\subsection*{Scope and status}
At $\Gamma=1$, the restriction entails full-vector conditional exchangeability; it is an assumption, not a conclusion obtained merely by adjusting for own covariates. For $\Gamma>1$, arbitrary point identification is not demanded. Full-vector positivity may deteriorate exponentially with cluster size and must enter any rate claim. This sensitivity-based subclass retains the attachment's concern about network confounding without silently declaring all confounders observed.
\subsection*{References for the source of the problem}
Carlos Cinelli, Avi Feller, Guido Imbens, Edward Kennedy, Sara Magliacane, and Jose Zubizarreta, \emph{Challenges in Statistics: A Dozen Challenges in Causality and Causal Inference}, arXiv:2508.17099v1, 23 August 2025, Section 3.2.3, paragraph \emph{Interference in observational studies}. \url{https://arxiv.org/html/2508.17099v1}.

\subsection*{Common conventions: Source mathematical conventions}

\textbf{Well-defined objects.}
All probability spaces and measurable variables are specified or inherited
from a local statistical model. Expectations used as finite targets require
the stated integrability. A decision rule or estimator is measurable with
respect to its available information; randomized rules may use an auxiliary
independent random variable. Norms without a subscript are Euclidean in
finite-dimensional spaces. An optimizer is requested only when attainment
is assured; otherwise the question concerns the optimal value and
$\varepsilon$-optimal admissible rules for every $\varepsilon>0$.

\textbf{Identification.}
A structural model $m\in\mathcal M$ implies an observable law
$P_m^{\rm obs}$ and target $\theta(m)$. The target is point identified on
$\mathcal M$ if
\[
 P_{m_1}^{\rm obs}=P_{m_2}^{\rm obs}
 \quad\Longrightarrow\quad\theta(m_1)=\theta(m_2)
 \qquad(m_1,m_2\in\mathcal M).
\]
When this fails, the sharp identified set is
\[
 \Theta(P;\mathcal M)=\{\theta(m):m\in\mathcal M,
                                  P_m^{\rm obs}=P\}.
\]
Specifying an intervention or estimand does not establish identification.
Positivity, exclusion, causal sufficiency, measurement restrictions, and
transport assumptions are substantive local inputs, never automatic
consequences of notation.

\textbf{Minimax risk and nontrivial inference.}
For a statistical experiment with data law $P^{(n)}$, target $\theta(P)$,
and nonnegative loss $\ell$, define
\[
 \mathcal R_n(\mathcal P,\ell)
   =\inf_{\widehat\theta_n}\sup_{P\in\mathcal P}
      \E_{P^{(n)}}\ell(\widehat\theta_n,\theta(P)).
\]
The infimum is over measurable estimators. A sharp rate is a sequence
$r_n$ for which $0<c\leq C<\infty$, independent of $n$, satisfies
$c r_n\leq\mathcal R_n\leq C r_n$ for all sufficiently large $n$.
Squared-error parametric risk is of order $n^{-1}$; the corresponding
root-mean-squared error is of order $n^{-1/2}$. These different conventions
are distinguished locally. A uniform asymptotic confidence region $C_n$
must satisfy
\[
 \liminf_{n\to\infty}\inf_{P\in\mathcal P}
     \Pr_{P^{(n)}}\{\theta(P)\in C_n\}\geq 1-\alpha,
 \qquad 0<\alpha<1,
\]
together with an informative diameter or expected-width requirement.
Coverage alone permits the vacuous whole parameter space. Testing questions
similarly impose both size control and power against a declared separated
alternative class.

\textbf{Binary causal baseline.}
When an entry explicitly invokes this baseline, one observes independent
copies of $O=(X,W,Y)$, with $W\in\{0,1\}$ and potential outcomes $Y(0),Y(1)$.
Assume consistency $Y=Y(W)$, no interference, conditional exchangeability
$(Y(0),Y(1))\perp W\mid X$, and overlap
$\epsilon\leq e(X)\leq1-\epsilon$ almost surely for a fixed
$0<\epsilon<1/2$, where
\[
 e(x)=\Pr(W=1\mid X=x),\qquad
 m_w(x)=\E[Y\mid W=w,X=x],\qquad
 \tau(x)=m_1(x)-m_0(x).
\]
These assumptions identify conditional mean effects almost everywhere
under the covariate law. Pointwise or uniform effects require appropriate
regularity and a specified support. Continuous treatment, longitudinal
data, adaptive experiments, and network interference require their own
assumptions in place of this baseline. Bounded outcomes, densities, and
smoothness are additional assumptions only where stated.

\textbf{Function classes.}
On $[0,1]^d$, $\mathcal H_d^s(L)$ is the isotropic Hölder ball of radius
$L>0$: put $k=\lceil s\rceil-1$ and $\gamma=s-k\in(0,1]$, bound derivatives
through order $k$, and require the order-$k$ derivatives to be
$\gamma$-Hölder with constant at most $L$. Thus integer orders use a
Lipschitz final derivative convention. The class
$\operatorname{BV}_d(L)$ consists of functions $f\in L^1((0,1)^d)$ whose
distributional derivative is a finite vector measure and for which
$\|f\|_{L^1}+|Df|((0,1)^d)\leq L$.
An $L^\infty$ bound or a particular pointwise version is imposed separately
when needed. Sparse and neural-network classes require a declared
dictionary or architecture and coefficient bounds.

An anisotropic H\"older class with declared block smoothness indices means
coordinatewise H\"older regularity in each block, uniformly in the other
blocks, with all corresponding derivative and H\"older bounds at most
the stated radius. Derivative orders and the integer-order convention in
each block are those just given. Mixed-derivative bounds are additional
restrictions only when explicitly stated.

\textbf{Decision and computational questions.}
Policy regret compares admissible feasible policies under the same law and
constraints. In computational entries, finite data and thresholds have
rational encodings; running time is measured in their total bit length.
Real-valued equilibrium search uses the explicitly declared algebraic or
FIXP encoding. An approximation ratio for maximization is written
$\operatorname{OPT}/\operatorname{ALG}$ and for minimization as
$\operatorname{ALG}/\operatorname{OPT}$, with zero optima and infeasible
instances handled locally. Historical approximation bounds are not
automatically current bounds.
% END ECONOMICS STATEMENT
\end{document}
